\chapter{Discussion and Conclusions}\label{ch:discussion}
The main demonstrated contribution is a reproducible Phase 1 substrate: fixed image splits and caption references, explicit provenance-oriented metadata, normalized visual and text embeddings, separate training-only indexes, and integrity checks. The retrieval smoke test shows that the stored image vectors can be searched and that the tested held-out queries do not retrieve themselves. This confirms pipeline operation under the recorded configuration.

Several boundaries matter for interpretation. The 0.8260 mean top-1 cosine is a similarity statistic without judged relevance labels. The successful QC status covers implemented checks, not the complete academic project. CLIP scene predictions are not calibrated; box relations are geometric; color and edge metrics are coarse image-level proxies. The selected 100-image validation set is small for broad statistical conclusions, and the 20-query smoke test is intended as a functional check rather than a benchmark. Full-scale retrieval relevance and generation quality remain unmeasured.

The two memory types create an opportunity for a controlled study. Image memory could retrieve visual neighbors; caption memory could retrieve language examples. A proper experiment must prevent the query's own references from entering its evidence, preserve train-only knowledge, keep generator settings fixed, and inspect whether retrieved examples help or propagate errors. The current separation of training and held-out data is a sound prerequisite, but prompt and generation leakage controls still need implementation.

\section{Limitations}\label{ch:limits}
The most important limitation is scope: Phase 2 is absent. There is no caption decoder, no RAG prompt assembly, no comparative method, no generation evaluation, no significance test, no ablation, and no hallucination study. There is also no original assignment PDF in the audited project locations, so the report cannot verify exact professor wording or the designated technique. The 100-versus-1,000 test-size conflict remains unresolved; the repository uses 1,000.

Reproducibility is partial. The seed, model identifier/revision, key tensor properties, and many artifact mappings are recorded, but a unified environment snapshot is missing. Historical counts of EXIF corrections and grayscale/RGBA conversions were not retained. The metadata audit flags raw-caption mapping integrity as partial even though five-per-image cardinality and current record counts pass. The preprocessing comparison was created for the audit and evaluates geometry only. These gaps should be addressed before the final comparative study.

Dataset limitations include the restricted 11,100-image sample and COCO's incomplete annotation coverage. Text token counts depend on tokenization. POS and zero-shot scene outputs may be wrong. Perceptual hash candidates indicate visual similarity rather than confirmed identity. These constraints limit claims to the observed artifacts and prevent strong generalization or causal statements.

\section{Conclusion and Next Steps}\label{ch:conclusion}
The repository contains a complete and quality-controlled Phase 1 preparation pipeline for a proposed RAG image-captioning study. It selects 10,000 training, 100 validation, and 1,000 test images, preserves five captions per image, prepares letterboxed RGB inputs, analyzes annotations and visual/text properties, extracts 512-dimensional normalized CLIP vectors, builds train-only image and caption indexes, and validates held-out image retrieval on 20 queries. Integrity checks report no split overlap, no exact duplicates, and no cross-split perceptual candidates at the configured threshold.

These findings establish data and retrieval readiness, not improved caption generation. The central comparative research question remains unanswered because generation experiments are pending. The next phase should obtain instructor clarification on test size and assigned technique; freeze an explicit generator and baseline; implement controlled image-, caption-, and hybrid-retrieval conditions; exclude query-associated references; evaluate with multiple complementary metrics and human grounding judgments; analyze paired outcomes and ablations; and publish complete configurations and outputs. Until then, all caption-quality claims must remain open.

\appendix
